\section{Introduction}
\label{sec:intro}
\begin{figure}[t]
\centering
\includegraphics[width=\columnwidth]{figures/fig1/fig1_overview.pdf}
\caption{The usual prompt sends four private reviews of this LaMP-3 user. The guard sends only
low-risk ones, with a cutoff tuned on 225 other users to pass at most half a private review on
average, and the answer is still 4 stars.}
\Description{Cartoon of one user. Top, on the laptop: the user asks to rate a new review that begins "Of all of the horse fly sprays, I like this the best". Below are six past reviews as numbered cards with short verbatim excerpts. Reviews 2 to 5 reveal private details, drawn in red and marked with a dog for pet or hearts for partner: "on my cocker", "on my Paso Fino", "MY COCKER", and "my poor, allergy sensitive girl". Reviews 1 and 6 are about weed killer and a cooling product and reveal nothing. Bottom, sent to the AI service: the usual way sends reviews 1 to 5, so four private reviews leave the laptop and the reader robot answers 4 stars. With the guard, only reviews 1 and 6 are sent, no private review leaves, and the answer is again 4 stars. The true rating is 4 stars.}
\label{fig:walkthrough}
\end{figure}
Personalized assistants answer better when the prompt carries records from
the user's history such as earlier posts, reviews or
messages~\cite{salemi2024lamp,salemi2025lampqa}. When the model runs at a
provider, each admitted record leaves the user's device, and personal
histories are full of self-disclosures about health, family, work and
location~\cite{dou2024reducing}. Data minimization asks that personal data be
limited to what the purpose needs~\cite{biega2020operationalizing,ganesh2025data},
which here means an on-device gate between the profile and the prompt.

Existing gates select profiles or personal fields for a task, reporting benchmark
utility and leakage with hand-set or utility-tuned thresholds~\cite{liu2026enough,du2026purple,zhou2026operationalizing,bagdasarian2024airgapagent,wang2026spmem,yu2026minim}.
Deployment instead needs an exposure bound at a chosen level for unseen users.

We obtain that number with conformal risk control~\cite{angelopoulos2024conformal}.
The exposure of an admission rule, the count of admitted records that an
auditor flags as self-disclosing, is bounded by the record cap and monotone
in the rule's one parameter, so calibrating that parameter on audited users
bounds the expected exposure of a new user's prompt by $\alpha$ for any
ranking of records (Figure~\ref{fig:walkthrough}). The statement is per
prompt, in expectation over users exchangeable with the calibration sample,
and relative to the auditor, which runs offline on consenting calibration
users only. The theory is standard; we study its use for record admission and its utility cost.

Across three tasks, two LLM auditors and three readers, we compare calibrated
admission with retrieval and a hindsight oracle. The experiments test the
bound's dependence on auditors and population, and measure the utility cost
of restricting whole records rather than removing disclosures within them.
